\documentclass[11pt]{article}

\usepackage[margin=1in]{geometry}
\usepackage{graphicx}
\usepackage{booktabs}
\usepackage{amsmath,amssymb}
\usepackage{float}
\usepackage[T1]{fontenc}
\usepackage{lmodern}
\usepackage{microtype}
\usepackage{caption}
\usepackage[numbers,square]{natbib}
\usepackage[hidelinks]{hyperref}
\captionsetup{font=small,labelfont=bf}

\graphicspath{{figures/}}

\title{\textbf{Optimizing the Rule Does Not Buy Cooperation:\\
A Null Result for Norm Optimization Across Six LLM Agents\\
in a Commons Dilemma}}
\author{}
\date{}

\begin{document}
\maketitle

\begin{abstract}
When self-interested agents share a renewable resource, a natural intervention is
to write them a better rule. We test whether \emph{automatically optimizing} that
rule improves cooperation, using two established textual optimizers---TextGrad and
Agentic Context Engineering (ACE)---applied to the shared norm in a five-agent
fishing commons. We run a $6\times3$ grid: six large language models playing the
agents, three optimizers (including a fixed-norm control), each for five
optimization iterations over two evaluation seeds, for $176$ complete $12$-month
episodes. The result is a clean null: in \emph{zero} of six models did either
optimizer produce a reward gain larger than the run-to-run variation of an
\emph{unchanged} norm. Pooled across models, all three arms---including the
control that cannot learn---trace the same iteration curve, confirming the
apparent ``improvements'' are measurement noise. What did move outcomes was the
\emph{agent}: mean reward ranged from $+2.42$ (Claude Opus~4.8, at the theoretical
ceiling) to $-0.18$ (GPT-5.6 Sol, which liquidated the fishery within one month
under every rule tried). The optimizers wrote visibly more sophisticated rules;
those rules simply did not score better. We read this as convergent evidence,
reached from the mechanism-design side, that in an unenforced setting the agent's
disposition---not the quality of the written rule---governs whether a commons
survives.
\end{abstract}

\section{Introduction}
The tragedy of the commons is the canonical failure of self-interested
coordination: when a renewable resource is shared, each actor's incentive is to
over-extract, and the resource collapses despite everyone preferring it to
survive~\citep{hardin1968,ostrom1990}. As large language model (LLM) agents are
placed in shared environments where they negotiate, allocate, and act with
minimal oversight, the same dilemma re-appears in software. GovSim~\citep{govsim}
established a concrete testbed: five LLM ``fishers'' draw from a lake that
regenerates each month, and the question is whether they can sustain it.

One intuitive lever is the \emph{rule}. If a community of agents can agree on a
written norm---``take a fair, sustainable share''---perhaps a better-written norm
produces better cooperation. This paper asks a sharp version of that question:
if we place a modern text optimizer in a loop, \emph{optimizing the norm against
measured outcomes}, does cooperation improve? Prompt- and text-optimization
methods such as TextGrad~\citep{textgrad} and ACE~\citep{ace} have shown large
gains on reasoning and agentic benchmarks; a commons norm is exactly the kind of
free-text ``variable'' they are designed to improve.

Our answer is negative, and cleanly so. Across six agent models and two
optimizers, no configuration produced an improvement distinguishable from the
noise of re-evaluating a \emph{fixed} norm. Meanwhile, swapping the model that
plays the agents moved outcomes across the entire available range. The rule was
close to inert; the agent was decisive.

This mirrors, from the opposite methodological direction, the central claim of
recent work on incomplete contracts among LLM agents~\citep{mechdesign}: that
designing an adequate mechanism is \emph{not} sufficient for cooperative
deployment, and that agents must be intrinsically prosocial. That work varies
prosociality and holds the mechanism family fixed; we hold disposition fixed
(all-selfish agents) and \emph{optimize the mechanism directly}. Both roads reach
the same place.

\paragraph{Contributions.}
\begin{enumerate}
  \item A $176$-episode, six-model, three-optimizer grid testing whether automated
  norm optimization improves cooperation in a commons dilemma.
  \item A clean null result, established against a fixed-norm control that supplies
  a per-model noise floor and exposes apparent gains as measurement noise.
  \item A large \emph{model} effect over the same episodes---from the theoretical
  optimum to total collapse---including a frontier model (GPT-5.6 Sol) that
  reproducibly destroys the fishery within one month.
  \item A methods analysis of \emph{why} the null arises and how a follow-up should
  be powered.
\end{enumerate}

\section{Setup}

\paragraph{Environment.}
We use the GovSim fishing commons~\citep{govsim}. A lake holds up to $100$ tons of
fish. Five agents each privately choose a monthly catch; catches execute
simultaneously and are made public. The surviving stock then regenerates by a
factor of $2$ (capped at $100$), so a full lake sustainably yields $50$ tons per
month---$10$ per agent. Sustained over-extraction shrinks next month's stock; a
lake driven to zero cannot recover and the episode ends. The horizon is $12$
months. Each agent is a separate instance of the same model, with its own private
memory; agents learn about each other only through the public harvest report and a
monthly negotiation phase.

\paragraph{The norm.}
Agents operate under a written natural-language norm (a ``constitution''). We use
the \emph{unenforced} regime: the norm is read by every agent but is not
mechanically imposed---each agent still freely chooses its catch, and the group may
amend or replace the norm each month if four of five agree. An unenforced norm is
the harder and more interesting test: it can only work if self-interested agents
find it worth \emph{following} and worth \emph{keeping}. All agents are given a
selfish disposition, isolating the norm as the intervention.

\paragraph{Reward.}
We score an episode with a single reward combining GovSim's metrics: total gain
$R$ (normalized against a fixed full-horizon benchmark so an early collapse cannot
score well), survival months $m$, equality $e = 1 - \text{Gini}$, and over-usage
$o$:
\[
  \text{reward} = \frac{R}{R_{\max}} + \frac{m}{12} + \tfrac{1}{2}\,e - o .
\]
Efficiency $u$ is reported but deliberately \emph{not} optimized: it saturates at
$1$ and an early collapse shrinks its own denominator, so maximizing it would
reward destroying the fishery---a failure mode we confirmed empirically (a
one-month collapse scores $u=1.00$). The reward's arithmetic maximum in this
setting is $2.5$.

\paragraph{Optimizers.}
Each optimizer proposes a norm, the norm is evaluated by running full episodes, and
the optimizer revises based on the measured reward. The three arms share this
interface and differ only in how the next norm is produced:
\begin{itemize}
  \item \textbf{Baseline (control).} The norm is held \emph{fixed} across all
  iterations. Its run-to-run variation is therefore pure evaluation noise---the
  yardstick every other arm must beat.
  \item \textbf{TextGrad}~\citep{textgrad}. An LLM writes a ``textual gradient''---a
  concrete critique of the current norm given its score---then applies it, with
  validation-based rollback of regressions.
  \item \textbf{ACE}~\citep{ace}. A Generator/Reflector/Curator loop grows a
  \emph{playbook} of transferable lessons and drafts norms from it.
\end{itemize}
To avoid confounding ``which model plays the commons better'' with ``which model
writes better norms,'' the optimizer's own LLM is held \emph{fixed} (Claude
Sonnet~5) across all cells; only the model playing the agents varies.

\paragraph{Grid.}
Six agent-model configurations, spanning size, vendor family, and reasoning:
Gemma~4~31B; GPT-5.6 Luna; Claude Sonnet~5 with reasoning disabled and enabled
(identical weights); Claude Opus~4.8; and GPT-5.6 Sol. Each cell runs all three
optimizers for five iterations over two evaluation seeds. Of $180$ planned
episodes, $176$ completed (four failed and are excluded; no cell lost more than
two). The grid cost \$444 in API usage and ran in $7.6$ hours.

\section{Results}

\subsection{The model dominates; the optimizer does not}
Figure~\ref{fig:reward} shows mean reward for every model$\times$optimizer cell.
The visual story is immediate: the bars are grouped tightly \emph{within} each
model and spread enormously \emph{across} models. Opus~4.8 sits at the ceiling
($2.5$); GPT-5.6 Sol sits below zero. Within any model, the three optimizers are
nearly indistinguishable.

\begin{figure}[H]
  \centering
  \includegraphics[width=\linewidth]{fig_reward_by_model.pdf}
  \caption{Mean reward by agent model and optimizer. The spread \emph{across}
  models (from the $2.5$ ceiling to below zero) dwarfs the spread \emph{across}
  optimizers within any model. Baseline is the fixed-norm control.}
  \label{fig:reward}
\end{figure}

Table~\ref{tab:main} gives the numbers. Ordered by mean reward, the six models
span $2.41$ (Opus) to $-0.18$ (Sol). The best-performing optimizer within each
model rarely matches the identity of the ``winner,'' and the differences are
small relative to the model effect.

\begin{table}[H]
  \centering
  \small
  \begin{tabular}{lcccc}
    \toprule
    Model & Baseline & TextGrad & ACE & Cell mean \\
    \midrule
    Opus 4.8            & $2.417$ & $2.462$ & $2.337$ & $\mathbf{2.405}$ \\
    Sonnet 5 (reasoning)& $2.146$ & $2.162$ & $2.251$ & $2.186$ \\
    Gemma 4 31B         & $2.033$ & $2.216$ & $2.263$ & $2.171$ \\
    Sonnet 5 (no reas.) & $2.087$ & $1.997$ & $2.224$ & $2.103$ \\
    GPT-5.6 Luna        & $1.764$ & $1.709$ & $0.926$ & $1.466$ \\
    GPT-5.6 Sol         & $-0.220$& $0.080$ & $-0.401$& $\mathbf{-0.180}$ \\
    \bottomrule
  \end{tabular}
  \caption{Mean reward over five iterations, by model and optimizer. The
  across-model range ($2.41 \rightarrow -0.18$) is an order of magnitude larger
  than any within-model optimizer difference.}
  \label{tab:main}
\end{table}

\subsection{The optimizer gains are noise}
Because the baseline arm holds the norm fixed, its spread across iterations
measures how loud the evaluation is. Figure~\ref{fig:noise} plots each optimizer's
best-over-baseline gain against that per-model noise floor. Every real gain falls
\emph{inside} the noise band. The single point that exceeds its floor---Sol under
TextGrad---is an artifact: its ``best'' occurred at iteration $0$, the
\emph{un-optimized} norm drawing a lucky episode, after which every iteration was
worse.

\begin{figure}[H]
  \centering
  \includegraphics[width=\linewidth]{fig_noise.pdf}
  \caption{Optimizer gain over the fixed-norm baseline (points) against the
  per-model noise floor (bars, the spread of the identical norm re-evaluated five
  times). Every genuine gain lies within the noise. The one point above its
  bar is an iteration-$0$ artifact, not an optimizer gain.}
  \label{fig:noise}
\end{figure}

The decisive diagnostic is Figure~\ref{fig:traj}. Pooled across the five
non-collapse models, all three arms---\emph{including the fixed-norm baseline,
which cannot learn}---trace the same trajectory: a rise to iteration $3$ and a
fall at iteration $4$, all ending below where they started. An identical shape in
an arm that is structurally incapable of improvement proves the shape is noise.
On a per-cell basis, TextGrad beats baseline in three of six cells and ACE in
three of six---a coin flip.

\begin{figure}[H]
  \centering
  \includegraphics[width=\linewidth]{fig_trajectories.pdf}
  \caption{Mean reward vs.\ optimizer iteration. \emph{Left:} pooled across the
  five non-collapse models; the fixed-norm baseline (which cannot learn) traces
  the same curve as the optimizers, so the curve is noise. Shaded bands are
  $\pm1$ SD across models. \emph{Right:} GPT-5.6 Sol, whose collapse regime no
  norm escapes.}
  \label{fig:traj}
\end{figure}

\subsection{The two frontier models are at opposite extremes}
Figure~\ref{fig:gainsurv} reports total gain and survival months per model.
Opus~4.8 essentially solves the commons---mean survival $12.0$ of $12$ months and
mean gain $548$, near the sustainable optimum of $600$. GPT-5.6 Sol does the
opposite: mean survival $1.6$ months and mean gain $130$. Inspection of Sol's
episodes shows two to three of its five agents demanding the \emph{entire} lake in
month one---collective demand reaching $6.4\times$ the sustainable level---across
every seed and under every norm the optimizers produced, including elaborate
stock-conditional quota systems with floor triggers and public reporting.
Sol's agents could find the sustainable action (two of five reliably requested
exactly $10$); others simultaneously went for everything.

\begin{figure}[H]
  \centering
  \includegraphics[width=\linewidth]{fig_gain_survival.pdf}
  \caption{Total gain $R$ (left) and survival months $m$ (right) by model, averaged
  over all episodes. Opus~4.8 approaches the sustainable optimum; GPT-5.6 Sol
  collapses the fishery in month one regardless of the norm in force.}
  \label{fig:gainsurv}
\end{figure}

Reasoning on versus off, tested on identical Sonnet~5 weights, produced
$2.186$ versus $2.103$---directionally favoring reasoning but well within the
noise floor. This grid cannot call it.

\subsection{The norms did evolve---into elaborate rules that did not help}
The null result is \emph{not} a failure of the optimizers to change the text.
Eleven of twelve optimizer runs found their best norm at a later iteration, and
the texts changed dramatically: from a single vague sentence to graduated,
state-contingent institutions. The starting norm was
\begin{quote}\itshape
Fishers should be mindful of the shared lake and avoid taking so much that the
fish cannot recover. Take a reasonable, fair share.
\end{quote}
The highest-scoring norm in the entire grid (Opus, TextGrad, reward $2.500$) was a
seven-clause rule with a public stock measurement, a piecewise catch schedule
conditioned on current stock, a total moratorium below the break-even point, an
equal-split requirement, a four-of-five amendment threshold, and an explicit
appeal to incentive compatibility. It out-scored the vague one-sentence norm by
$0.022$, against a noise floor of $0.167$.

The interpretation we defend: for a capable model, an elaborate norm is worth
almost nothing over a vague one---the model supplies the cooperation and the rule
is nearly decorative---while for an uncooperative model an elaborate norm is
\emph{also} worth nothing, because it defects regardless. The norm text had
leverage in neither regime.

\section{Discussion}

\paragraph{A null with teeth.}
The design could not have detected a moderate optimizer effect, and this is the
most important caveat. With two evaluation seeds, the noise floor reached $1.43$
reward points on a $2.5$-point scale; the effect we sought is plausibly $0.1$--$0.2$.
The null is therefore ``not detected,'' not ``does not exist.'' Three changes
would power a follow-up: (i) more seeds per evaluation---noise on a mean falls as
$1/\sqrt{n}$, so $2\rightarrow8$ halves it; (ii) restricting optimization to
mid-range models with headroom (Opus begins $0.02$ from the ceiling and cannot
demonstrate improvement); and (iii) harder regeneration regimes (stochastic or
hysteretic), which lower baseline scores and open the room a deterministic
setting closes. Retaining the fixed-norm control is essential: without it, the
iteration-$3$ peaks would have been reported as evidence that optimization works.

\paragraph{Relation to mechanism design.}
\citet{mechdesign} prove that when contracts cannot distinguish all future
contingencies there is a welfare loss no mechanism can eliminate, and that
prosocial agents close it. They vary disposition and hold the mechanism family
fixed. We do the complement: hold disposition fixed and optimize the mechanism
text directly, with two methods designed for exactly this kind of optimization.
Optimizing the rule bought nothing measurable; changing the agent bought
everything. Designing an adequate mechanism was not sufficient.

\paragraph{Limits.}
We test one condition---unenforced natural-language norms, all-selfish agents,
deterministic regeneration. Enforced (code) contracts, where a well-written rule
has mechanical teeth, are exactly where norm quality should matter most and were
not tested. The optimizer LLM was fixed by design, so we did not test whether a
model optimizes best for itself. A hardcoded three-month planning horizon in the
agent prompt affects all cells equally and is a plausible contributor to Sol's
month-one collapse. All null results are at $n=2$ seeds.

\section{Conclusion}
Across six LLM agents and two established text optimizers, automatically optimizing
a commons norm produced no cooperation gain distinguishable from the noise of
re-evaluating a fixed norm, while the choice of agent model swung outcomes from
the theoretical optimum to total collapse over the very same episodes. In an
unenforced setting, the agent---not the rule---decides whether the commons
survives. The rule can be made elaborate; it cannot be made to matter.

\begin{thebibliography}{9}
\small
\bibitem[Hardin(1968)]{hardin1968}
G.~Hardin. The Tragedy of the Commons. \emph{Science}, 162(3859):1243--1248, 1968.
\bibitem[Ostrom(1990)]{ostrom1990}
E.~Ostrom. \emph{Governing the Commons}. Cambridge University Press, 1990.
\bibitem[Piatti et al.(2024)]{govsim}
G.~Piatti, Z.~Jin, M.~Kleiman-Weiner, B.~Sch\"olkopf, M.~Sachan, and R.~Mihalcea.
Cooperate or Collapse: Emergence of Sustainable Cooperation in a Society of LLM
Agents. 2024.
\bibitem[Huang et al.(2026)]{mechdesign}
X.~A.~Huang, C.~Tharas, S.~Marro, V.~Q.~Truong, B.~Sch\"olkopf, E.~La~Malfa, and
Z.~Jin. Mechanism Design Is Not Enough: Prosocial Agents for Cooperative AI. 2026.
\bibitem[Yuksekgonul et al.(2025)]{textgrad}
M.~Yuksekgonul et al. TextGrad: Automatic ``Differentiation'' via Text.
\emph{Nature}, 2025.
\bibitem[ACE(2025)]{ace}
Agentic Context Engineering: Evolving Contexts for Self-Improving Language Models.
2025.
\end{thebibliography}

\end{document}
